\documentclass[11pt]{article}
\usepackage[margin=1in]{geometry}
\usepackage[slantedGreek]{mathpazo}
\usepackage[pdftex]{graphicx}
\usepackage{amsmath,amssymb,amsfonts}
\usepackage{booktabs}
\usepackage{float}
\usepackage{caption}
\usepackage[numbers,sort&compress]{natbib}
\usepackage{hyperref}
\hypersetup{colorlinks=true,linkcolor=blue,citecolor=blue,urlcolor=blue}

\newcommand{\sandtable}{\textsc{SandTable}}
% placeholder until tools/make_bench_numbers.py runs against the landed benchmark
\newcommand{\BnTruthReps}{400}

% generated by tools/make_results_numbers.py; do not edit
\newcommand{\RsCells}{3{,}600}
\newcommand{\RsRows}{36{,}000}
\newcommand{\RsReplicates}{20}
\newcommand{\RsMethods}{10}
\newcommand{\RsProblems}{30}
\newcommand{\RsGradeable}{13}
\newcommand{\RsRegretOnly}{17}
\newcommand{\RsReachSep}{14}
\newcommand{\RsLowGpEi}{0.365}
\newcommand{\RsLowGpNoisy}{0.354}
\newcommand{\RsLowGpRaw}{0.364}
\newcommand{\RsLowGrid}{0.220}
\newcommand{\RsLowHalving}{0.294}
\newcommand{\RsLowRandom}{0.290}
\newcommand{\RsLowThompson}{0.341}
\newcommand{\RsLowUcb}{0.344}
\newcommand{\RsLowUcbPooled}{0.335}
\newcommand{\RsLowUcbRaw}{0.346}
\newcommand{\RsLowWinner}{grid}
\newcommand{\RsLowWinnerRegret}{0.220}
\newcommand{\RsMidGpEi}{0.149}
\newcommand{\RsMidGpNoisy}{0.171}
\newcommand{\RsMidGpRaw}{0.110}
\newcommand{\RsMidGrid}{0.093}
\newcommand{\RsMidHalving}{0.093}
\newcommand{\RsMidRandom}{0.115}
\newcommand{\RsMidThompson}{0.104}
\newcommand{\RsMidUcb}{0.128}
\newcommand{\RsMidUcbPooled}{0.107}
\newcommand{\RsMidUcbRaw}{0.105}
\newcommand{\RsMidWinner}{halving}
\newcommand{\RsMidWinnerRegret}{0.093}
\newcommand{\RsTopGpEi}{0.135}
\newcommand{\RsTopGpNoisy}{0.148}
\newcommand{\RsTopGpRaw}{0.057}
\newcommand{\RsTopGrid}{0.058}
\newcommand{\RsTopHalving}{0.061}
\newcommand{\RsTopRandom}{0.078}
\newcommand{\RsTopThompson}{0.048}
\newcommand{\RsTopUcb}{0.048}
\newcommand{\RsTopUcbPooled}{0.040}
\newcommand{\RsTopUcbRaw}{0.054}
\newcommand{\RsTopWinner}{ucb\_pooled}
\newcommand{\RsTopWinnerRegret}{0.040}
\newcommand{\RsGpEiVsGridEff}{+0.0764}
\newcommand{\RsGpEiVsGridT}{+8.14}
\newcommand{\RsGpEiVsGridVerdict}{gate up}
\newcommand{\RsGpNoisyVsGridEff}{+0.0894}
\newcommand{\RsGpNoisyVsGridT}{+9.36}
\newcommand{\RsGpNoisyVsGridVerdict}{gate up}
\newcommand{\RsGpRawVsGridEff}{-0.0010}
\newcommand{\RsGpRawVsGridT}{-0.22}
\newcommand{\RsGpRawVsGridVerdict}{equivalent}
\newcommand{\RsHalvingVsGridEff}{+0.0024}
\newcommand{\RsHalvingVsGridT}{+0.53}
\newcommand{\RsHalvingVsGridVerdict}{equivalent}
\newcommand{\RsRandomVsGridEff}{+0.0191}
\newcommand{\RsRandomVsGridT}{+4.27}
\newcommand{\RsRandomVsGridVerdict}{gate up}
\newcommand{\RsThompsonVsGridEff}{-0.0107}
\newcommand{\RsThompsonVsGridT}{-1.99}
\newcommand{\RsThompsonVsGridVerdict}{equivalent}
\newcommand{\RsUcbVsGridEff}{-0.0100}
\newcommand{\RsUcbVsGridT}{-1.90}
\newcommand{\RsUcbVsGridVerdict}{equivalent}
\newcommand{\RsUcbPooledVsGridEff}{-0.0183}
\newcommand{\RsUcbPooledVsGridT}{-3.74}
\newcommand{\RsUcbPooledVsGridVerdict}{gate down}
\newcommand{\RsUcbRawVsGridEff}{-0.0049}
\newcommand{\RsUcbRawVsGridT}{-1.62}
\newcommand{\RsUcbRawVsGridVerdict}{equivalent}
\newcommand{\RsGpEiVsRandomEff}{+0.0573}
\newcommand{\RsGpEiVsRandomT}{+5.68}
\newcommand{\RsGpEiVsRandomVerdict}{gate up}
\newcommand{\RsGpNoisyVsRandomEff}{+0.0703}
\newcommand{\RsGpNoisyVsRandomT}{+6.85}
\newcommand{\RsGpNoisyVsRandomVerdict}{gate up}
\newcommand{\RsGpRawVsRandomEff}{-0.0201}
\newcommand{\RsGpRawVsRandomT}{-3.49}
\newcommand{\RsGpRawVsRandomVerdict}{gate down}
\newcommand{\RsGridVsRandomEff}{-0.0191}
\newcommand{\RsGridVsRandomT}{-4.27}
\newcommand{\RsGridVsRandomVerdict}{gate down}
\newcommand{\RsHalvingVsRandomEff}{-0.0166}
\newcommand{\RsHalvingVsRandomT}{-3.05}
\newcommand{\RsHalvingVsRandomVerdict}{gate down}
\newcommand{\RsThompsonVsRandomEff}{-0.0298}
\newcommand{\RsThompsonVsRandomT}{-4.59}
\newcommand{\RsThompsonVsRandomVerdict}{gate down}
\newcommand{\RsUcbVsRandomEff}{-0.0291}
\newcommand{\RsUcbVsRandomT}{-4.52}
\newcommand{\RsUcbVsRandomVerdict}{gate down}
\newcommand{\RsUcbPooledVsRandomEff}{-0.0374}
\newcommand{\RsUcbPooledVsRandomT}{-6.22}
\newcommand{\RsUcbPooledVsRandomVerdict}{gate down}
\newcommand{\RsUcbRawVsRandomEff}{-0.0240}
\newcommand{\RsUcbRawVsRandomT}{-5.05}
\newcommand{\RsUcbRawVsRandomVerdict}{gate down}
\newcommand{\RsBeatGrid}{1}
\newcommand{\RsBeatGridNames}{ucb\_pooled}
\newcommand{\RsBeatRandom}{7}
\newcommand{\RsBeatRandomNames}{gp\_raw, grid, halving, thompson, ucb, ucb\_pooled, ucb\_raw}
\newcommand{\RsFoundLowGpEi}{25}
\newcommand{\RsFoundLowGpNoisy}{25}
\newcommand{\RsFoundLowGpRaw}{25}
\newcommand{\RsFoundLowGrid}{73}
\newcommand{\RsFoundLowHalving}{68}
\newcommand{\RsFoundLowRandom}{53}
\newcommand{\RsFoundLowThompson}{25}
\newcommand{\RsFoundLowUcb}{25}
\newcommand{\RsFoundLowUcbPooled}{25}
\newcommand{\RsFoundLowUcbRaw}{25}
\newcommand{\RsFoundLowWinner}{grid}
\newcommand{\RsFoundMidGpEi}{52}
\newcommand{\RsFoundMidGpNoisy}{49}
\newcommand{\RsFoundMidGpRaw}{64}
\newcommand{\RsFoundMidGrid}{82}
\newcommand{\RsFoundMidHalving}{83}
\newcommand{\RsFoundMidRandom}{64}
\newcommand{\RsFoundMidThompson}{66}
\newcommand{\RsFoundMidUcb}{67}
\newcommand{\RsFoundMidUcbPooled}{67}
\newcommand{\RsFoundMidUcbRaw}{66}
\newcommand{\RsFoundMidWinner}{halving}
\newcommand{\RsFoundTopGpEi}{63}
\newcommand{\RsFoundTopGpNoisy}{60}
\newcommand{\RsFoundTopGpRaw}{91}
\newcommand{\RsFoundTopGrid}{91}
\newcommand{\RsFoundTopHalving}{89}
\newcommand{\RsFoundTopRandom}{75}
\newcommand{\RsFoundTopThompson}{89}
\newcommand{\RsFoundTopUcb}{88}
\newcommand{\RsFoundTopUcbPooled}{92}
\newcommand{\RsFoundTopUcbRaw}{90}
\newcommand{\RsFoundTopWinner}{ucb\_pooled}
\newcommand{\RsCrossLow}{grid}
\newcommand{\RsCrossTop}{ucb\_pooled}
\newcommand{\RsCrossLowEff}{+0.1147}
\newcommand{\RsCrossLowT}{+8.23}
\newcommand{\RsCrossLowSign}{75}
\newcommand{\RsCrossLowVerdict}{gate up}
\newcommand{\RsCrossLowN}{600}
\newcommand{\RsCrossTopEff}{-0.0183}
\newcommand{\RsCrossTopT}{-3.74}
\newcommand{\RsCrossTopSign}{71}
\newcommand{\RsCrossTopVerdict}{gate down}
\newcommand{\RsCrossTopN}{600}
\newcommand{\RsCrossLowUni}{0.220}
\newcommand{\RsCrossLowUniFound}{73}
\newcommand{\RsCrossLowAdapt}{0.335}
\newcommand{\RsCrossLowAdaptFound}{25}
\newcommand{\RsCrossMidUni}{0.093}
\newcommand{\RsCrossMidUniFound}{82}
\newcommand{\RsCrossMidAdapt}{0.107}
\newcommand{\RsCrossMidAdaptFound}{67}
\newcommand{\RsCrossTopUni}{0.058}
\newcommand{\RsCrossTopUniFound}{91}
\newcommand{\RsCrossTopAdapt}{0.040}
\newcommand{\RsCrossTopAdaptFound}{92}
\newcommand{\RsCrossTopBinds}{clears both clauses of the tier}
\newcommand{\RsCrossReal}{yes}
\newcommand{\RsDiscVarEff}{-0.0194}
\newcommand{\RsDiscVarT}{-5.94}
\newcommand{\RsDiscVarSign}{63}
\newcommand{\RsDiscVarVerdict}{gate down}
\newcommand{\RsDiscVarN}{3{,}600}
\newcommand{\RsDiscIncEff}{+0.0106}
\newcommand{\RsDiscIncT}{+2.81}
\newcommand{\RsDiscIncSign}{56}
\newcommand{\RsDiscIncVerdict}{equivalent}
\newcommand{\RsDiscIncN}{3{,}600}
\newcommand{\RsDiscScaleEff}{+0.0082}
\newcommand{\RsDiscScaleT}{+2.52}
\newcommand{\RsDiscScaleSign}{54}
\newcommand{\RsDiscScaleVerdict}{equivalent}
\newcommand{\RsDiscScaleN}{3{,}600}
\newcommand{\RsDiscPriorEff}{+0.0343}
\newcommand{\RsDiscPriorT}{+10.13}
\newcommand{\RsDiscPriorSign}{69}
\newcommand{\RsDiscPriorVerdict}{gate up}
\newcommand{\RsDiscPriorN}{3{,}600}
\newcommand{\RsReuse}{5{,}642}
\newcommand{\RsReuseRows}{4.0}


\title{Searching a mission design space: what a fast simulator changes about\\
       optimization, and what it does not}
\author{}
\date{}

\begin{document}
\maketitle

\begin{abstract}
Mission design asks which setting of a simulator's controllable parameters gives the best outcome,
and the standard answer is Bayesian optimization. We test that answer on a suite of
\BnKeyProblems{} mission-design search problems whose optima were established by separate registered
studies of the simulator that involved no optimizer, with the budget denominated in mission
replications rather than evaluations and every method paired on the same missions in the same order.
Across \RsRows{} graded runs at \RsReplicates{} replicates per cell, \RsBeatGrid{} of the nine
adaptive methods beat uniform allocation over the design set at any budget on the decision tier we
registered before running them. Below the replication at which a design set can be ranked at all,
uniform wins outright, leading the strongest adaptive allocator by \RsCrossLowEff{} of the
achievable spread at $t = \RsCrossLowT{}$ and recovering the best design \RsCrossLowUniFound{}\% of
the time against \RsCrossLowAdaptFound{}\%. Above it a variance-pooled bandit does pull ahead,
\RsCrossTopAdapt{} normalized regret against \RsCrossTopUni{}, but the paired contrast
\RsCrossTopBinds{}: the difference is real and reproducible at $t = \RsCrossTopT{}$ and it is
smaller than the \RsCrossTopMpe{} of spread we had registered as worth acting on. The cause is measured rather than
argued, in a column the sweep records for every run: at the budgets where uniform wins, the adaptive
methods evaluate \RsLowCoverAdapt{}\% of the design set and uniform evaluates \RsLowCoverUni{}\%. An
allocator can only concentrate on evidence that ranks the arms, on these problems that evidence is
expensive, and until it is affordable the allocator is spending its budget on a fraction of the
designs. A design that is never evaluated cannot be returned. This is the third answer the
same comparison has given us. The first, uniform unbeatable everywhere, was an artefact of a budget
ladder whose top rung sat below the replication at which anything is learnable; the second, a
crossing that reversed the ranking, was an artefact of \RsPrevReplicates{} replicates. Both are reported here
with the sweeps that produced them, because a benchmark that reports only its final answer gives a
reader no way to tell which kind of artefact it is still carrying.

Establishing the answer key made the cost measurable, and it is severe. Only \BnKeyResolved{} of the
\BnKeyProblems{} problems separate their own top two designs at two standard errors, none of the
mission-success problems do, and on that axis a key built from \BnRepsSucc{} missions per design can
only separate two designs differing by a median \BnKeyDetSucc{} percent of the entire distance from
best design to worst. The spread across the design set is nonetheless large, at least
\BnSpreadOverGapMed{} times the gap between the best two, so choosing a good design is affordable
while proving which is best is not, by a factor of at least \BnCheaper{}. We therefore grade
normalized regret and carry argmax recovery as a flag, and we argue that a requirement written from
the argmax of a mission sweep is usually written from noise.
\end{abstract}

\section{Introduction}
\label{sec:intro}

A mission simulator that runs in a fifth of a second changes what a design study can be. When one
run took a day, the design was chosen by judgement and the simulator was used to defend it. When a
run is cheap enough that a laptop can do ten thousand overnight, the design can be searched for, and
the question becomes which search to run. The obvious answer is Bayesian optimization, which is what
one reaches for when evaluations are expensive and the space is continuous \citep{shahriari2016taking,
brochu2010tutorial}. Mission design is neither. The evaluations are cheap and the space is small,
discrete and often unordered, and both of those facts point somewhere else.

They point at a harder problem hiding inside the easy one. A mission simulator does not return the
value of a design, it returns a draw from a distribution over missions, so an evaluation has a
precision that the searcher chooses and pays for. Looking at ten designs once and one design ten
times cost the same and answer different questions. That makes mission design a best-arm
identification problem before it is an optimization problem \citep{bubeck2011pure,karnin2013almost},
and it makes the usual benchmarking convention, counting evaluations, quietly wrong: a method that
buys more replications per evaluation looks better without being better.

There is a second thing the literature holds fixed that this setting does not. An optimizer benchmark
varies the method and holds the objective constant, because on a test function the objective is the
problem. A mission has no single objective. The same missions can be scored on delivered service, on
whether the mission succeeded, or on what fraction of the force was lost, and a program office cares
about all three. We began this work expecting the interesting variation to lie between methods. It
does not. On identical missions the three axes differ by more than an order of magnitude in the
replications needed to tell the best design from the runner-up, and, more awkwardly, the ordering
among them is not stable: the axis that is cheapest to estimate at one operating point is the
dearest a few assets away. A benchmark that varies only the method cannot see any of that, and a
design study that picks its metric for reasons of reporting convention has made a larger decision
than it knows.

\paragraph{Contributions} We give a benchmark for mission-design search whose answer key is not
self-reported, and use it to argue that the axis a study searches on constrains its result more
tightly than the optimizer it searches with.

\begin{itemize}
\item A suite of mission-design search problems on a fast multi-agent simulator, six of them
      carrying an optimum established by a separate registered study that involved no optimizer, and
      the rest by a high-precision table computed on seeds disjoint from every search
      (Section~\ref{sec:bench}).
\item The measurement that most of those problems have no affordable best design at all. Only
      \BnKeyResolved{} of \BnKeyProblems{} separate their own top two designs at two standard
      errors, \BnKeyResSucc{} of \BnKeyNSucc{} on mission success, and on that axis the smallest
      gap a key can detect is a median \BnKeyDetSucc{}\% of the entire distance from best design to
      worst. We check first that this is a cost and not a bound effect, since a metric pinned at
      its ceiling would tie its top two as a matter of definition (Section~\ref{sec:bench}).
\item The consequence for what a benchmark should grade. The spread from best design to worst is a
      median \BnSpreadOverGapMed{} times the gap between the best two, so distinguishing the
      returned design from the best one is about \BnCheaper{} times cheaper than distinguishing the
      best from the runner-up. Choosing a design well is affordable and proving which is best often
      is not, so the suite grades normalized regret and carries argmax recovery as a flag
      (Section~\ref{sec:bench}).
\item A budget denominated in mission replications rather than evaluations, with the per-problem
      ladder derived from a measured quantity rather than chosen, and capped by the number of real
      missions standing behind the design it is spent on (Sections~\ref{sec:bench} and
      \ref{sec:methods}).
\item A comparison of \BnMethods{} methods under that budget, paired exactly rather than nominally:
      the answer key's per-mission readings are kept, so replicate $r$ fixes one mission ordering per
      design and every method that reaches a design receives the same missions in the same order.
      Two of the methods carry a deliberately unrepaired twin and two carry a variant differing in a
      single estimator, so a broken estimator can be told apart from a claim about adaptive search
      (Sections~\ref{sec:methods} and \ref{sec:results}).
\item The finding that the cost of estimating a mission metric is a property of the pair (metric,
      operating point) rather than of the metric, established by holding one design set fixed and
      scoring the same missions on three axes at four operating points (Section~\ref{sec:results}).
\item The consequence for method design, which appears inside a method rather than between two of
      them: an exploration constant is denominated in the metric's noise, because the textbook
      bandit bonus \citep{auer2002finite}, derived for rewards on the unit interval, exceeds the gap it must resolve by
      \BnBonusOverGapWorst{} times here. We kept the unscaled version as a control rather than
      repairing it silently, and the control is what showed our reasoning for the repair to be
      wrong: we registered that the unscaled allocator would behave as uniform sampling and it does
      not, it is \RsMOneWorstDir{} uniform allocation at the low rung by \RsMOneWorstMag{} on
      \RsMOneWorstField{} (Sections~\ref{sec:methods} and \ref{sec:results}).
\end{itemize}

The honest summary of the result is a negative one about method choice and a positive one about
study design, and Section~\ref{sec:discussion} argues that this is the more useful pair for a
program office than a ranking of optimizers would have been.

\section{The design-space search problem}
\label{sec:problem}

A mission design is a setting of the parameters a program office can actually choose: how many
assets one operator supervises, how deep a console lets its backlog grow, which request it works
next, how long a rebuffed vehicle waits before asking again, how far apart a column drives. A
mission simulator turns a design into a distribution over outcomes, and a design decision needs the
best design in a space of them. Written that way it is an optimization problem, and the standard
answer is Bayesian optimization: fit a surrogate to what has been evaluated, and let an acquisition
function choose where to look next.

Two properties of the mission-design setting make that answer less obvious than it looks, and both
of them are properties a fast simulator creates rather than removes.

The first is that the objective is a Monte-Carlo average, so every evaluation has a precision the
searcher chooses and pays for. Spending a thousand missions on one design and ten on another are
both legal moves, and a search method's real decision is jointly which design to try and how hard
to look at it. A benchmark that counts evaluations rather than replications makes that decision
invisible and quietly rewards whichever method buys the most precision per evaluation.

The second is that mission design spaces are usually small, often discrete, and frequently
categorical. A queue discipline is one of seven names with no ordering among them. A backlog bound
is one integer. A control modality is one of three. On a space of that shape the question "where
should I look next" is close to answered by enumeration, and the question that remains, "which of
these arms is best given noisy draws", is the one the bandit literature was built for
\citep{bubeck2011pure,karnin2013almost}. Best-arm identification and optimization are different
problems with different optimal strategies, and treating the second as the first is a choice that
should be tested rather than assumed.

\section{The benchmark, and why its answer key is not self-reported}
\label{sec:bench}

Optimizer benchmarks usually score a method against the best value that method or its competitors
happened to find. That is a self-report: it can say which method did better on a run, and it cannot
say whether any of them found the right design, because the right design is not independently
known. On synthetic test functions the answer is known and the problem is that the function is not a
mission. This benchmark takes the other route.

Every problem in the suite is a design space of the mission simulator \sandtable{}, and every
optimum in it was established by a separate study of that simulator with no view of any search. Each
of those studies registered a conjecture and a falsifier before its run, gated its contrasts as
paired blockwise comparisons at a fixed tier, and reported retractions where the falsifier fired.
None of them involved an optimizer. So "did the method find the best design" is a question with an
answer that was written down by somebody else, first.

The suite is deliberately not a set of variations on one shape, because the claim under test is that
the shape decides which family of methods wins.

\begin{itemize}
\item \textbf{One integer, two fleets, two answers.} Sizing the console's backlog: seven capacities,
      optimum at two on a fleet whose vehicles all want the same thing and at one on a fleet whose
      deadlines differ. The smallest problem in the suite, and the one where enumeration is
      genuinely competitive, which is why it is here. A method that cannot match a grid on seven
      designs has no business on a larger space.
\item \textbf{A categorical axis with no ordering.} Seven queue disciplines, named rather than
      numbered. This is where a surrogate over a one-hot encoding is weakest and a bandit allocator
      is most natural, and it is common in mission design: a control modality, a doctrine, a
      formation type.
\item \textbf{A two-dimensional mixed space.} The backlog bound crossed with the span it is sized
      for, twenty-eight designs, the first problem where enumeration is expensive enough to be worth
      beating.
\item \textbf{A pair that is identical to the searcher and has different answers.} How long a
      rebuffed vehicle should wait before asking again has an interior optimum on a console that
      cannot refuse work and no optimum at all on a console that can. The two problems present the
      same variable over the same domain, and any method that returns one answer to the question
      ``what should this parameter be'' is wrong on one of them. It is in the suite as a trap, and
      it is a trap real design studies fall into.
\end{itemize}

\paragraph{Auditing the key.} An answer key can be wrong, so it is checked rather than trusted.
Every design in every problem is evaluated on a seed block disjoint from the ones the searches draw
on, and the resulting table is compared against the registered optimum. Agreement is reported per
problem and disagreement is reported rather than reconciled: a table that put the optimum somewhere
else would be a finding about the original study or about the precision of the table, and hiding it
would defeat the purpose of having a key at all.

The replication count behind the key is per metric and not one number, which is the paper's own
argument applied to its own instrument. \BnRepsSrv{} replications resolve the top two designs on
delivered service; the loss fraction needs \BnRepsLoss{}, a factor of \BnRepsRatio{}. A flat count
set for the first axis would have bought a precise key on the cheap axis and a noisy one on the
dear axis, and the noisy key would then have graded search methods on exactly the axis this paper
claims is expensive to estimate. The key would have been least trustworthy precisely where the
paper looks hardest.

Two further conventions matter more than they sound. Every design draws the SAME mission seeds, so
the comparison between designs is paired and what has to be resolved is the variance of the
difference rather than the variance of each arm; an unpaired key needs several times the
replications for the same resolving power and the noisiest axes pay the most. And every table
records the precision it actually achieved, rather than assuming the replication count bought the
precision it was chosen for.

\paragraph{Most of the benchmark has no best design, and that is a result rather than a defect.}
Recording achieved precision turns the question ``how expensive is this metric'' into one that can
be answered instead of estimated, and the answer is severe. Of the \BnKeyProblems{} problems, only
\BnKeyResolved{} separate their own top two designs at two standard errors, at replication counts
of \BnRepsSrv{} to \BnRepsLoss{} missions per design. By axis: \BnKeyResSrv{} of \BnKeyNSrv{} on
delivered service, \BnKeyResLoss{} of \BnKeyNLoss{} on blue loss, and \BnKeyResSucc{} of
\BnKeyNSucc{} on mission success.

The natural next number is how many missions an unresolved table would have needed, and we do not
report one, because it cannot be measured from a table that failed. That extrapolation divides by
the observed gap, and a table is called unresolved precisely when its gap sits inside two standard
errors of zero, so the quantity is ill-conditioned exactly on the problems it would be quoted for.
It is close to entailed as well: a gap of zero returns the replication count already spent and a
gap at the resolving threshold returns a quarter of it, so whatever the data say the answer lands
in a fixed band. On \BnKeyNoisy{} of the \BnKeyProblems{} tables the observed gap is smaller than
one standard error of itself, and on the worst of them the extrapolation read fifty-two million
missions per design from a gap of nine parts in one hundred thousand measured to plus or minus
thirty-eight hundredths of a percent.

What can be measured is the resolving power itself, against the room the design axis has, and both
halves of that are estimated directly. The smallest top-two gap a key can detect is two standard
errors of the difference; the room is the distance from best design to worst. On delivered service
the key detects a gap of \BnKeyDetSrv{}\% of that room, on blue loss \BnKeyDetLoss{}\%, and on
mission success \BnKeyDetSucc{}\%. On success, in other words, an answer key that spent
\BnRepsSucc{} missions per design can only separate two designs that differ by about half of
everything the design axis is able to do. \BnKeyUnrankable{} problem is worse than that:
\BnKeyDetWorstSuccKey{} cannot detect a gap smaller than \BnKeyDetWorstSucc{}\% of its own spread,
so it cannot rank its design set at all. That is what ``no affordable answer'' means here, stated
without dividing by a quantity consistent with zero.

One check has to come before that is read as a cost. A metric pinned at its bound would tie its top
two designs as a matter of definition, and no replication would ever separate them; the tie would be
a theorem and reporting it as expensive would be wrong. It is not the explanation here, and the test is reachability rather
than eyeballing: on \BnBoundMoves{} of the \BnBoundNSucc{} success problems the design axis moves
mission success further than that problem's own answer key can resolve, by a factor of at least
\BnBoundMovesMin{}, so a tie between the top two designs is a statement about the designs and not
one the bound entails. The widest of them, \BnBoundWideKey{}, spans \BnBoundWideSpan{} from best
design to worst with the best at \BnBoundWideBest{} and \BnBoundWideHead{} of room still above it.
The single exception is \BnBoundCeilKey{}, whose best design reads \BnBoundCeilBest{} across a span
of \BnBoundCeilSpan{} against a detectable gap of \BnBoundCeilDet{}: there the axis moves the metric
less than the key can see, and it is reported as ceiling-adjacent rather than as a replication cost.

\paragraph{Choosing well and proving which is best differ by two orders of magnitude.} The
structure underneath the previous paragraph is the useful part, and it is not that these problems
are flat. The design axis moves the metric a great deal; it is the top two designs that are close.
Across the suite the spread from best design to worst is at least a median
\BnSpreadOverGapMed{} times the gap between the best two, and at least \BnSpreadOverGapMax{} at the
extreme. Both are floors rather than estimates: on the \BnSpreadBounded{} problems whose gap is not
resolved the gap is known only to be at most two standard errors, and dividing by that ceiling
gives a bound from two measured quantities, while on the rest the same expression returns the ratio
exactly. Since the replication needed to separate two designs goes as the square of that ratio,
distinguishing the returned design from the best one is at least \BnCheaper{} times cheaper than
distinguishing the best from the runner-up. Choosing a design well is worth a lot and is affordable. Proving which single design is
best is worth little and frequently is not.

That is an argument about what to grade, and the benchmark follows it. A problem whose key cannot
separate its own top two is kept and scored on regret rather than dropped, with the flag recorded on
every row so either grading can be recovered from one sweep. Dropping them would have discarded
\BnKeyUnresolved{} of \BnKeyProblems{} problems, including every problem on the metric a sponsor
actually reads, on the grounds that a question nobody needs answered could not be answered. The
sweep nonetheless grades argmax on one more problem than the tables resolve, because
\BnArgmaxFromStudy{} of the unresolved problems carries an optimum established by a separate
registered study and so has a known argmax its own table cannot recover. Both counts are reported
because they are different questions, and a reader who meets only one of them will find the other
inexplicable.

\paragraph{Two scores, because they disagree.} A method is scored on whether it returned the best
design, which is a yes or a no, and on how much mission performance the design it returned gives up
against the best one, which is a number. Regret is carried as a fraction of the achievable spread,
zero for the best design and one for the worst, so that problems scored on delivered service,
mission success and blue loss can be averaged at all; in raw metric units they cannot be. These two
scores do not agree, and the disagreement is one of the paper's results rather than a nuisance. On
the backlog-sizing problem the best two capacities differ by under one percent of delivered service,
so a method that returns the runner-up has almost no regret and has still got the answer wrong.
Which failure matters depends on what the search is for. A study choosing a setting cares about
regret; a study writing a requirement cares about the answer. The measurement above says which of
those two questions a benchmark can afford to grade.

\section{Methods under a common budget}
\label{sec:methods}

\BnMethods{} methods are compared, and the comparison is made fair by two conventions that are
stated here because they are where most of the honesty lives.

\paragraph{The budget is missions, not evaluations.} Every method is given a fixed number of mission
replications and may spend them however it likes: many designs looked at briefly, few designs looked
at hard, or anything between. A method that asks for more than it has left is stopped and recorded
as having failed to return inside its budget, which is a property of the method and not a fault in
the test. Counting evaluations instead would let a method buy precision for free, and precision is
exactly what a stochastic objective makes expensive. It would also hide the decision that
distinguishes these methods most: with a noisy objective, choosing how hard to look at a design is
as consequential as choosing which design to look at.

\paragraph{Methods are paired on the mission randomness, exactly rather than nominally.} Replicate
$r$ fixes one mission ordering per design, and a method that reaches a design receives the same
missions in the same order as every other method on that replicate. So a difference between two
methods within a replicate is a difference in what they chose to spend that stream on and not in
what they happened to draw. This is the common random numbers discipline the simulator's own
statistical gate uses, applied one level up, and it is worth saying why it takes more than seeding
the methods alike. How a method spends its budget is the thing under measurement, so two methods
handed the same generator consume it differently and diverge after their first disagreement. The
pairing is made exact instead by drawing the comparison from stored per-mission readings: the
answer key's replications are kept rather than summarized, and a pull of $n$ missions on a design
takes the next $n$ of that design's ordering. What is resampled is the simulator's own mission
distribution, at \BnRepsLoss{} real missions per design on the dearest axis, and no ladder is
allowed a rung deeper than the pool behind it, so no method is ever served the same mission twice
while appearing to buy fresh precision.

\paragraph{Three arms exist to separate a defect from a finding.} Two of the methods below carry a
deliberately unrepaired twin and two carry a variant that differs in one estimator. A scaled
allocator that loses to its unscaled control on the noisy axes could mean that adaptive allocation
fails there, or it could mean that a variance estimated from two observations collapsed and starved
an arm; those two readings are indistinguishable from the results alone, and are told apart by an
arm that pools the estimate across designs and changes nothing else. The same holds for the
surrogate, where an incumbent taken as the best observed value is biased upward by noise. The
broken versions are kept rather than quietly fixed, because a benchmark that reports only repaired
methods reports a selection.

The methods span three families.

\emph{Enumeration.} A grid spends the budget evenly over every design and returns the best sample
mean. Random search draws designs uniformly with a fixed batch each. Both are here because they are
what a program office actually does and because they are the baselines that benchmarks most often
omit and most often lose to.

\emph{Bandit allocation.} Successive halving cuts the field in half each round and re-spends the
freed budget on the survivors \citep{karnin2013almost,jamieson2015nonstochastic}. An upper-confidence
allocator and a Thompson sampler spend a fixed chunk of missions per decision and choose the next
design from what they have seen. None of the three fits a response surface or needs the design space
to have a geometry, which is what makes them usable on a categorical axis. The objective they are
built for is pure exploration, identifying the best arm, which is a different objective from
minimizing cumulative regret and has different optimal strategies \citep{bubeck2011pure}. Mission
design wants the first.

\emph{Bayesian surrogate.} A Gaussian process over the encoded design space with an
expected-improvement acquisition \citep{brochu2010tutorial,shahriari2016taking}, ordinal axes scaled
to the unit interval and categorical axes one-hot. It is implemented directly over the enumerated
design set rather than through an optimization package. On spaces this small the acquisition can be
scored exhaustively, so the part a package spends most of its effort on is not the part in question;
what is in question is whether the surrogate's smoothing helps when the axis carries no order or the
response is a cliff, and a package's defaults would hide that.

\paragraph{What is not being claimed.} The surrogate here is a plain one. A better-tuned kernel, a
learned lengthscale, or a heteroscedastic noise model would all improve it, and none of them would
change the argument, which is about what the shape of a mission design space does to the choice of
method rather than about which implementation is best. Where the surrogate loses to a grid on seven
designs, the conclusion drawn is that the problem did not need a surrogate, not that surrogates do
not work.

\paragraph{Priority, and an older literature that has it.} Framing this as budget allocation across
competing designs whose values are known only through noisy replication is not new, and the
literature that owns it is simulation optimization rather than machine learning. Optimal computing
budget allocation asks precisely this question and answers it for the normal-means case
\citep{chen2000simulation}. The transfer being made here is in the other direction from the usual
one: not a new method, but a demonstration that a mission-design space has the shape that
literature assumed and not the shape the surrogate literature assumes.

\section{Results}
\label{sec:results}


The sweep is \RsRows{} graded runs: \RsMethods{} methods on \RsProblems{} problems at six budget
rungs each, \RsReplicates{} replicates per cell, every method on every replicate drawing the same
missions in the same order from the stored pools. \RsGradeable{} of the problems can be graded on
argmax recovery and all of them on normalized regret.

At \RsReplicates{} replicates every paired contrast below carries tens of thousands of matched
runs, enough that all of them reach a $t$ statistic above ten. That matters for how the verdicts
should be read: where a contrast is reported as equivalent, it is not that the comparison was
underpowered, it is that the difference is smaller than the tier's minimum practical effect. An
earlier \RsPrevReplicates{}-replicate version of this same sweep returned two verdicts the full run
does not support, and both are reported below rather than quietly replaced.

\paragraph{Uniform allocation is not beaten at any budget, and at small budgets it wins by a wide
margin.} Means are taken over the \RsCells{} problem, budget and replicate cells that every one of
the \RsMethods{} methods reached, so no two are averaged over different denominators.

At the bottom rung the lowest mean normalized regret belongs to \RsCrossLow{}, at
\RsCrossLowUni{} of the achievable spread against \RsCrossLowAdapt{} for the strongest adaptive
allocator. Paired on the replicate across \RsCrossLowN{} matched runs that is \RsCrossLowEff{} in
favour of uniform allocation at $t = \RsCrossLowT{}$ with \RsCrossLowSign{}\% sign consistency.
Argmax recovery reads the same way and more starkly, \RsCrossLowUniFound{}\% against
\RsCrossLowAdaptFound{}\%.

At the top rung the mean ordering changes places, \RsCrossTopAdapt{} for \RsCrossTop{} against
\RsCrossTopUni{}, and the change does not clear the tier: the paired effect is \RsCrossTopEff{} and
\RsCrossTopBinds{}, which here is \RsCrossTopMpe{} in units of the contrast's own spread. It is not
a question of precision. At \RsReplicates{} replicates the contrast reaches $t = \RsCrossTopT{}$
with \RsCrossTopSign{}\% sign consistency; it is simply too small a difference to matter. Across the
whole method set, \RsBeatGrid{} of the nine others separate from uniform allocation in the
favourable direction at the top of the ladder and \RsWorseGrid{} separate in the unfavourable one:
\RsWorseGridNames{}. Both counts are given, because a table reporting only the first would say that
nothing separates from uniform allocation about a set in which the two surrogate methods lose to it
by \RsGpEiVsGridEff{} and \RsGpNoisyVsGridEff{} of the spread. Argmax recovery has converged by
then, \RsCrossTopUniFound{}\% against \RsCrossTopAdaptFound{}\%. Table~\ref{tab:methods} gives every
method at all three rungs, and it does not say what we expected it to. The gap between best and
worst method barely narrows with budget, \RsRangeLow{} of the spread at the bottom rung against
\RsRangeTop{} at the top. What changes is which methods are outliers. The interquartile range falls
from \RsIqrLow{} to \RsIqrTop{}, so the middle of the field converges tightly while the range is
held open by the two expected-improvement surrogates, which end at \RsTopGpEi{} and \RsTopGpNoisy{}
against \RsCrossTopAdapt{} for the best method. At a budget that can rank the design set, the
methods that work are indistinguishable from each other and from enumeration, and the methods that
do not work fail by a wide margin. That is a benchmark separating working from broken, which is a
weaker thing than separating good from better, and it is what this suite can support.

\begin{table}[H]
\centering
\caption{Every method at three points on each problem's budget ladder, averaged over the
\RsCells{} cells all \RsMethods{} methods reached. Normalized regret is the fraction of the
achievable spread the returned design gives up, so lower is better and zero is the best design.
Argmax recovery is over the \RsGradeable{} problems whose key can rank its design set. Rows are
ordered by top-rung regret; the two enumeration baselines are italicised. The ordering at the top
rung is real and is not actionable: the best adaptive method leads uniform allocation by less than
the \RsCrossTopMpe{} of spread the tier asks for.}
\label{tab:methods}
% generated by tools/make_results_numbers.py; do not edit
\begin{tabular}{lrrrrrr}
\toprule
& \multicolumn{3}{c}{normalized regret} & \multicolumn{3}{c}{argmax recovery (\%)} \\
\cmidrule(lr){2-4}\cmidrule(lr){5-7}
method & low & mid & top & low & mid & top \\
\midrule
ucb\_pooled & 0.310 & 0.102 & 0.041 & 29 & 71 & 92 \\
thompson & 0.310 & 0.106 & 0.046 & 29 & 70 & 90 \\
ucb & 0.322 & 0.125 & 0.049 & 29 & 69 & 89 \\
ucb\_raw & 0.314 & 0.112 & 0.050 & 29 & 69 & 91 \\
\textit{grid} & 0.204 & 0.099 & 0.051 & 69 & 82 & 91 \\
gp\_raw & 0.346 & 0.114 & 0.051 & 29 & 69 & 91 \\
halving & 0.293 & 0.105 & 0.057 & 65 & 79 & 89 \\
\textit{random} & 0.271 & 0.121 & 0.067 & 56 & 67 & 75 \\
gp\_ei & 0.347 & 0.156 & 0.137 & 29 & 56 & 68 \\
gp\_noisy & 0.343 & 0.179 & 0.157 & 29 & 54 & 64 \\
\bottomrule
\end{tabular}

\end{table}

\paragraph{This is the third answer this comparison has given, and the first two were artefacts of
the study design.} The sequence is reported because each step was caused by something a reader
should check in their own benchmark. The first sweep found uniform allocation unbeatable at every
budget, and that was a property of the ladder: sizing the top rung to four times the budget that
separates best design from worst put it below the budget that separates the top TWO on
\RsOldLadderShort{} of the \RsOldLadderHave{} problems, in one case by a factor of a thousand, so
every method was allocating on evidence that cannot order the arms and uniform allocation is
trivially optimal there. The corrected ladder reaches the separating budget wherever the stored
mission pools allow, which is \RsReachSep{} of \RsLadderHave{} problems; on the remaining
\RsLadderShort{} the top rung is still short, and the finding below is therefore a statement about
the regime this suite can afford to reach rather than about every budget a study might have. With the ladder corrected to reach the separating budget, a
\RsPrevReplicates{}-replicate sweep found a crossing: the same pair of methods gave
\RsPrevCrossTopEff{} at the top rung at $t = \RsPrevCrossTopT{}$, which the tier read as
\RsPrevCrossTopVerdict{}, adaptive allocation in front. At \RsReplicates{} replicates that crossing
does not survive. The effect falls to \RsCrossTopEff{} and lands under the minimum practical
effect, while the $t$ statistic more than DOUBLES in magnitude, from
$\RsPrevCrossTopT{}$ to $\RsCrossTopT{}$, because the fifteenfold increase in matched runs outruns
the shrinking estimate. That combination is the useful part: more
data made the difference more certain and less worth acting on at the same time, and a tier written
only in significance would have taken the wrong lesson from the extra replicates. The finding is
one-directional after all, but for a different reason than the first sweep gave, and only the third
version rests on a ladder that spans the regime and a replicate count that settles the contrast.

\paragraph{Why adaptive allocation is a liability at small budgets and no help at large ones.} An
adaptive allocator earns its keep by spending less on designs it has ruled out, and ruling a design
out means separating it from the leader. Section~\ref{sec:bench} measured what separation costs on
these problems. Below that budget a searcher is deciding where to allocate on evidence that cannot
rank the arms it is choosing between, so it concentrates on noise. Every row of the sweep records
how many DISTINCT designs its method actually evaluated, and at the bottom rung that column settles
the question. On problems averaging \RsLowDesigns{} designs, \RsCrossTop{} evaluates
\RsLowTouchAdapt{} of them, \RsLowCoverAdapt{}\% of the set, while \RsCrossLow{} evaluates
\RsLowTouchUni{}, \RsLowCoverUni{}\%. The adaptive methods are not choosing badly among the designs
they have seen. They are never seeing most of the design set, and a design that was never evaluated
cannot be returned. That is why the deficit at the bottom rung is wide rather than marginal, and why
it shows up far more sharply in argmax recovery than in regret: committing early to two or three
designs still lands somewhere reasonable in regret terms while picking the wrong design outright.
Above that budget the evidence is finally worth acting on, and by then there is nothing left to win,
because on a design space this small uniform allocation has already reached every arm to a
precision that ranks them.

This gives a program office a rule it can apply before choosing a method, using a quantity the
answer key already reports. Compare the study's per-design budget against the replication at which
the design set can be ranked. Below it, enumerate; the adaptive machinery will cost real
performance, and it will cost it mostly in the answer rather than in the outcome. Above it, the
choice stops mattering on a space of this size, so enumerate anyway and spend the effort on the
answer key instead.

\paragraph{Only one of four estimator choices makes a difference that matters.} Four of the
\RsMethods{} methods exist as members of a pair differing in exactly one estimator, so that a reader
can tell an implementation choice apart from a claim about adaptive search. Each pair is gated on
\RsDiscVarN{} matched runs, which is enough that every one of the four reaches a $t$ statistic
above ten. Three of them are still equivalent at the tier, because the tier asks for a difference
that is large as well as consistent.

Denominating the exploration bonus in the metric's own noise, rather than in the unit interval the
textbook derivation assumes \citep{auer2002finite}, does not repay itself: \RsDiscScaleEff{} at $t = \RsDiscScaleT{}$ with
\RsDiscScaleSign{}\% sign consistency, which the tier reads as \RsDiscScaleVerdict{}. Taking the
surrogate's incumbent from the posterior mean rather than from the best observed value is
\RsDiscIncVerdict{} at \RsDiscIncEff{}. Pooling the per-mission standard deviation across arms
rather than estimating it per arm is \RsDiscVarVerdict{} at \RsDiscVarEff{} despite $t =
\RsDiscVarT{}$; at \RsPrevReplicates{}
replicates this one gated, and it is the clearest case in the paper of a difference that is real,
reproducible, and too small to act on.

The one estimator choice that does clear the tier runs against the method it was supposed to
improve. Fitting expected improvement to the best observed value rather than to the raw surrogate
is worse, by \RsDiscPriorEff{} at $t = \RsDiscPriorT{}$ with \RsDiscPriorSign{}\% sign
consistency, the one contrast of the four the tier scores as \RsDiscPriorVerdict{}. An incumbent read off the noisiest arm the searcher has yet seen is an
optimistic target, and chasing it costs more than the acquisition rule recovers. That is worth
saying plainly, because three of these four differences are exactly the kind of tuning a comparison
without controls would have reported as a method contribution, and the fourth would have been
reported with its sign the wrong way round.

\paragraph{Registered verdicts, reported whichever way they went.} Three of the five predictions
registered before the sweep failed. CR-72 M1 held that the unscaled allocator would be equivalent
to uniform allocation at every budget on both grading rules, since a bonus two orders of magnitude
larger than the gap it must resolve should collapse it into round-robin sampling. It is
\RsMOneVerdict{}: \RsMOneGates{} of the \RsMOneOf{} rung and grading-rule combinations gate away
from uniform allocation, worst at the \RsMOneWorstRung{} rung on \RsMOneWorstField{}, where the
unscaled allocator is \RsMOneWorstDir{} uniform allocation by \RsMOneWorstMag{} at $t =
\RsMOneWorstT{}$. The registration was right that the bonus swamps the gap and wrong that this
makes the allocator uniform, and the coverage column above says why. A bonus that always selects
the least-pulled arm still has to pull it in batches, and at a budget of a few missions per design
the batches run out before the arms do: the unscaled allocator reaches \RsLowCoverRaw{}\% of the
design set at the bottom rung where uniform allocation reaches \RsLowCoverUni{}\%. Degenerating
toward round-robin is not the same as completing one, and the difference is exactly the designs it
never evaluates. At the top rung, where the budget covers every design either way, the two ARE
equivalent, so scoring M1 there alone would have called it held. CR-72 M2 held that scaling the bonus would recover the loss at the middle rung. It is
\RsMTwoVerdict{} at \RsMTwoEff{}, $t = \RsMTwoT{}$, \RsMTwoSign{}\% sign consistency. CR-72 M3 cannot be scored at all, because it asked for an
ordering across three axes and the mission-success axis has no problem whose key can rank its
design set, which Section~\ref{sec:bench} had already measured and the registration had not
allowed for.

CR-73 B held that argmax recovery and regret would rank methods differently, so that the choice of
grading rule mattered. It is retracted: the two rankings agree at Spearman
\RsRankSpearman{} and the same method wins both on \RsRankSame{} of the \RsRankProblems{}
problems, which clears both branches of its falsifier: the
registration asked for agreement above nine tenths OR the same winner on ten of twelve problems,
and both fired. The grading rule still matters for which problems can be graded,
which is the claim Section~\ref{sec:bench} makes and this one does not touch, but the ranking
argument was academic and is withdrawn.

CR-73 C holds, and by a narrower margin than the registration expected. Its falsifier said the
\RsSepUnresN{} problems the sweep cannot grade on argmax should be dropped if NO method separates
from \texttt{random} under regret. \RsSepUnresEither{} of \RsSepUnresOf{} do, in opposite
directions: \RsSepUnresRegNames{} is better than random and \RsSepUnresWorseNames{} is worse, both
at the registered two clauses. The remaining \RsSepUnresTied{} are indistinguishable from random on
these problems, almost always because sign consistency falls short of sixty percent rather than because
the effect is small. So the problems do carry information about method quality, which is what
justified keeping them rather than dropping more than half the suite, but what they mainly separate
is a bad method from an arbitrary one. A benchmark built only from them would be a weak instrument,
and the paper does not claim otherwise.

\paragraph{Pool reuse.} A method that asks for more missions of a design than the pool holds
resamples draws it has already seen, and the precision it reports is then optimistic. Across the
whole sweep that happened \RsReuse{} times, on \RsReuseRows{}\% of runs, and every affected row
carries the count.

\section{Discussion}
\label{sec:discussion}

The result a program office should take from this is not that Bayesian optimization is a poor
method. It is that on a mission-design problem the binding constraint is the cost of estimating an
outcome, not the cleverness of the search over designs. A design space of seven queue capacities is
enumerable; what is expensive is knowing which of the seven is better, and that expense is set by
the metric and the operating point rather than by the optimizer. Every method in this comparison spends the same
missions, and which way of spending them wins is decided by whether the budget clears the cost of
ranking the arms. The literature that motivates surrogate search assumes the opposite regime, a
space too large to enumerate and an evaluation whose value is known once taken, and neither half of
that assumption holds below the ranking threshold.

That reading is what makes the result portable rather than a property of this simulator. The
condition under which uniform allocation wins is stated in quantities a study can measure before
choosing a method: the design set is small enough to enumerate inside the budget, and the budget per
design is below the replication at which the set can be ranked at all. Both are reported by the
answer key, so a study can tell which side of that threshold it is on without running the comparison
at all. Below it, uniform is not merely safe but strictly better and the choice is settled. Above
it, the best adaptive method is ahead by less than the tier calls actionable, so the choice is
settled the other way, by whichever method is cheaper to defend. Where a mission study has a genuinely large or continuous design space, or an outcome
cheap enough that its ranking is settled early, the argument does not apply and the surrogate
literature should be expected to hold.

The methodological half may matter more than the comparison. An optimizer benchmark that scores
methods against the best value any of them found cannot tell whether a method found the right
design, and on a stochastic simulator it cannot tell whether the difference it reports is larger
than the noise it was measured in. Establishing the key separately, denominating the budget in
mission replications, and pairing methods on the same missions are three changes that together turn
a comparison of self-reports into a measurement. The third of them is cheap and worth stating
plainly: storing the answer key's own per-mission readings and resampling them made pairing exact
rather than nominal, at a cost of a few megabytes, and it removed the dominant source of variance
in every contrast in this paper.

The measurement that most of the suite has no affordable best design deserves separating from the
comparison, because it survives whatever the comparison says. It is the reason this benchmark grades
normalized regret. It is also a caution about the wider practice of establishing a requirement by
running a sweep and reading off the argmax, which is a common way for a simulator to influence a
programmatic decision. On \BnKeyNSucc{} mission-success problems here that argmax is not
identifiable at \BnRepsSucc{} missions per design, and the spread between best and worst design is
real and large while the gap between the top two is not measurable. A requirement written from the
argmax of such a sweep is written from noise; a requirement written from the achievable spread is
not.

The limits are worth stating in the same terms. This is one simulator, two design knobs, three
outcome metrics, and design spaces of \BnDesignsMin{} to \BnDesignsMax{} points, all discrete. The
methods were given a fixed budget and no tuning, which is the right convention for a comparison and
the wrong one for a deployment. The comparison rests on \RsReplicates{} replicates per cell, enough
that the contrasts reported as equivalent are equivalent on effect size and not merely unresolved.

The sharpest limit is the one the ladder still imposes. Even after the correction, the top rung
reaches the budget that separates the top two designs on only \RsReachSep{} of \RsLadderHave{}
problems, because the stored mission pools cap it on the rest. So the claim that uniform allocation
is not beaten holds over the budgets this suite can afford to buy, not over every budget a study
might have. What would settle it is deeper pools on the problems that are short, which is a
measurement cost rather than a method question, and we say what it would take rather than leaving
the reader to find the gap.

And the search operates on a static design space; the question of whether a learned policy can do
better by choosing where to look as a sequential decision problem, rather than by fitting a
surrogate to what it has seen, is the natural next step and is not answered here.

\bibliographystyle{plainnat}
\bibliography{ref}

\end{document}
